\section{分母の微分が分子にあるとき}\label{節:対数型の置換}

$\frac1t$ の原始関数は $\log|t|$ だから，置換積分を使うと
\[
 \int\frac{g'(x)}{g(x)}\dd x=\log|g(x)|+C
\]
となる。ただし，$g(x)\ne0$ である区間で使う。$g$ が連続ならば，その区間内で符号を変えるには $0$ を通る必要があるから，符号は一定である。正なら $\log g(x)$，負なら $\log(-g(x))$ を微分すれば，同じ商に戻る。

\begin{example}{絶対値を残すべき対数}
$\int\frac{6x}{3x^2-2}\dd x$ を求めよ。
\end{example}

分母の $3x^2-2$ を微分すると，分子の $6x$ が現れる。$t=3x^2-2$ と置けば $\dd t=6x\dd x$ なので，
\[
 \int\frac{6x}{3x^2-2}\dd x
 =\int\frac1t\dd t
 =\log|t|+C
 =\log|3x^2-2|+C.
\]
分母が $0$ になる $x=\pm\sqrt{\frac23}$ は除く。この二点の間では $3x^2-2<0$ なので，絶対値を外して $\log(3x^2-2)$ とすると実数の式ではなくなる。一方，$\log|3x^2-2|$ は三つの区間のどこでも使える形である。任意定数は各区間で独立に選んでよい。

分母が二次式だから必ず部分分数にする，という決め方はしない。この例では分子を見れば分母の微分があり，一度の置換で終わる。前章の式の整理と置換積分は競合する規則ではなく，式の特徴に応じて選ぶ道具である。

ただし，$\int\frac1{3x^2-2}\dd x$ の分子は $6x$ ではないから，同じ公式をそのまま使うことはできない。分母だけを $t$ と換えて，$\dd x$ を変えずに残すと，変数の異なる式を混ぜたことになる。置換後は，積分される式と微分の記号の両方を，新しい変数で書き直す必要がある。

